\chapter{Central Banking in Practice}\label{ch:b3:central-banking}

A central bank sets one interest rate, overnight, and wants to move inflation two
years ahead. Between the two lie its balance sheet, the plumbing of the money
market, the yield curve and the expectations of everyone who prices a loan. This
chapter follows the chain: how the rate is set in practice, how guidance and asset
purchases reach longer rates, how economists measure the effect of a policy
surprise --- with one estimate made in the chapter --- what limits on mortgages do,
and what the inflation of 2021--2023 taught.

\begin{openingchart}{Central-bank total assets as a percent of nominal GDP, year
ends 2003--2025: the Federal Reserve and the Eurosystem.}{Board of Governors of the
Federal Reserve System, H.4.1 (RESPPMA\_N.WW); European Central Bank, consolidated
balance sheet of the Eurosystem; World Bank, World Development Indicators; Eurostat
(nama\_10\_gdp); ratios are ours}
\begin{tikzpicture}
\begin{axis}[omtime, width=0.7\linewidth, height=0.42\linewidth, xmin=2002, xtick distance=5, xmax=2025.5,
  ymin=0, ymax=75, ylabel={Percent of GDP}, legend pos=north west,
  legend style={font=\footnotesize}]
  \addplot[omDef, thick, mark=none] table[x=year, y=fed] {figdata/bachelor-3/15-central-banking/out/cb_assets.csv};
  \addplot[omThm, thick, mark=none] table[x=year, y=ecb] {figdata/bachelor-3/15-central-banking/out/cb_assets.csv};
  \legend{Federal Reserve, Eurosystem}
\end{axis}
\end{tikzpicture}
\end{openingchart}

% figdata: cb_assets_points
At the end of 2007 the Federal Reserve's assets were 6~\% of U.S. GDP; by 2014 they were
26~\%, by 2021 37~\%, and in 2025 still 22~\%. The Eurosystem's went from 16~\% of
euro-area GDP to 68~\% in 2021, and were 39~\% in 2025. In fifteen years the two
largest central banks became, in size, the largest investors in their own bond
markets.

\begin{recall}
The grades 10--12 volume explained how banks and the central bank create money.
The Year 2 volume treated inflation targets, credibility and independence, the
policy rule and the zero lower bound, and the lender of last resort.
\end{recall}

\begin{notation}[Symbols of this chapter]\label{not:b3:central-banking:symbols}
$i$ is the policy rate, $i_D$ and $i_L$ the deposit and lending rates of the central
bank, $r$ the overnight market rate, $R$ \omterm{def:b3:central-banking:balance-sheet}{bank reserves} at the central bank, $y_n$ the
yield on an $n$-year bond and $\mathit{tp}_n$ its \omterm{def:b3:central-banking:qe}{term premium}.
\end{notation}

\section{The balance sheet}\label{sec:b3:central-banking:balance-sheet}

\begin{definition}[Central-bank balance sheet; bank reserves]\label{def:b3:central-banking:balance-sheet}
\index{central-bank balance sheet}\index{bank reserves}
A \emph{central bank's balance sheet} records on the asset side the securities it
holds, its loans to banks, gold and foreign reserves, and on the liability side the
currency in circulation, the government's deposits and \emph{bank reserves}: the
deposits of commercial banks at the central bank, which they use to settle
payments with each other.
\end{definition}

\begin{measure}[What a central bank's balance sheet records]
When the central bank buys a bond from a bank, it pays by crediting the bank's
reserve account: assets and reserves rise by the same amount, and no one else can
create or destroy reserves. Banks can pass reserves among themselves but not,
together, get rid of them --- except by withdrawing currency, or when the
government's account at the central bank grows. The total of reserves is the central
bank's choice; what banks do with them is not.
\end{measure}

\section{Setting the overnight rate}\label{sec:b3:central-banking:corridor}

\begin{definition}[Corridor system; floor system; standing facility]\label{def:b3:central-banking:corridor}
\index{corridor system}\index{floor system}\index{standing facility}
A \emph{standing facility} is a central-bank offer open to any eligible bank at a
fixed rate: a deposit facility at $i_D$ and a lending facility at $i_L > i_D$. In a
\emph{corridor system} the central bank supplies scarce reserves and steers the
overnight rate between the two. In a \emph{floor system} it supplies abundant
reserves, and the overnight rate sits at the deposit rate.
\end{definition}

\begin{proposition}[The corridor and the floor]\label{prop:b3:central-banking:corridor}
The overnight rate satisfies $i_D \le r \le i_L$. If banks' demand for reserves is
satiated at $\bar R$ and reserves exceed $\bar R$, then $r = i_D$.
\end{proposition}

\begin{proof}
No bank lends overnight below $i_D$, because it can deposit at the central bank at
$i_D$; no bank borrows above $i_L$, because it can borrow from the central bank at
$i_L$. If reserves exceed what banks want to hold at any rate above $i_D$, each bank
with extra reserves tries to lend them; competition drives the rate down to the
point where holding them is as good as lending, $i_D$.
\end{proof}

\begin{omfigure}
\begin{tikzpicture}[econ, x=0.09cm, y=1.6cm]
  \draw[econ axes] (0,0) -- (95,0) node[right, font=\scriptsize] {reserves};
  \draw[econ axes] (0,0) -- (0,3.4) node[above, font=\scriptsize] {overnight rate};
  \draw[guide] (0,2.5) node[left, font=\scriptsize] {$i_L$} -- (90,2.5);
  \draw[guide] (0,1.5) node[left, font=\scriptsize] {$i_D$} -- (90,1.5);
  \draw[curve, omDef] (5,2.5) -- (20,2.5) -- (60,1.5) -- (90,1.5)
    node[above left, font=\scriptsize] {demand};
  \draw[curve, omThm] (40,0) -- (40,3.1) node[above, font=\scriptsize] {corridor};
  \draw[curve, omThm, dashed] (80,0) -- (80,3.1) node[above, font=\scriptsize] {floor};
  \fill (40,2) circle (2pt) node[right, font=\scriptsize] {$r = i$};
  \fill (80,1.5) circle (2pt);
  \foreach \v in {20,40,60,80} \draw (\v,0) -- (\v,-0.06) node[below, font=\scriptsize] {\v};
\end{tikzpicture}
\par\smallskip{\small The reserve market of the weekend problem. Demand slopes between
20 (below which banks pay the lending rate) and 60 (above which they are satiated).
With 40 of reserves the overnight rate is mid-corridor; with 80 it sits on the
floor.}
\end{omfigure}

A central bank that runs a corridor must forecast reserve demand and adjust the
supply every day to keep the overnight rate near its target. Once asset purchases
have made reserves abundant, it sets the overnight rate by setting the rate it
pays on reserves: the policy rate and the size of the balance sheet become two
separate instruments.

\section{Guidance and the yield curve}\label{sec:b3:central-banking:guidance}

\begin{proposition}[The long rate: expectations plus a term premium]\label{prop:b3:central-banking:term}
The $n$-year yield can be written $y_n = \frac1n\sum_{k=0}^{n-1}E_ti_{t+k} + \mathit{tp}_n$: the average expected
short rate over the bond's life plus a \omterm{def:b3:central-banking:qe}{term premium}.
\end{proposition}

\begin{proof}
An investor can hold an $n$-year bond or roll over one-year bonds. If the two paid
the same expected return, the long yield would equal the average expected short
rate (the expectations hypothesis, to first order in logs). Whatever difference
remains is, by definition, the \omterm{def:b3:central-banking:qe}{term premium}: the extra return investors require
for bearing the risk of the long bond.
\end{proof}

\begin{definition}[Forward guidance]\label{def:b3:central-banking:forward-guidance}
\index{forward guidance}
\emph{Forward guidance} is a central bank's communication about the future path of
its policy rate, meant to move the expected short rates in the long yield.
\end{definition}

The proposition names the two levers left when the short rate is at its floor:
lower the expected path (guidance), or lower the \omterm{def:b3:central-banking:qe}{term premium} (asset purchases).
A promise to keep rates low after recovery is worth what markets believe: the
Year 2 volume showed that such a promise is not time-consistent, so its effect
depends on the bank's credibility, and markets read it in the futures curve.

\begin{omchart}{The federal funds rate and the ten-year Treasury yield, monthly averages,
2005--2025 (percent).}{Board of Governors of the Federal Reserve System, H.15
(RIFSPFF\_N.M, RIFLGFCY10\_N.M)}
\begin{tikzpicture}
\begin{axis}[omtime, width=0.7\linewidth, height=0.4\linewidth, xmin=2004, xtick distance=5, xmax=2026.2,
  ymin=0, ymax=6, ylabel={Percent}, legend columns=2,
  legend style={at={(0.5,1.03)}, anchor=south, font=\footnotesize}]
  \addplot[omDef, thick, mark=none] table[x=t, y=ff] {figdata/bachelor-3/15-central-banking/out/rates.csv};
  \addplot[omThm, thick, mark=none] table[x=t, y=y10] {figdata/bachelor-3/15-central-banking/out/rates.csv};
  \legend{Federal funds, Ten-year Treasury}
\end{axis}
\end{tikzpicture}
\end{omchart}

\section{Asset purchases}\label{sec:b3:central-banking:qe}

\begin{definition}[Quantitative easing; term premium]\label{def:b3:central-banking:qe}
\index{quantitative easing}\index{term premium}
\emph{Quantitative easing} (large-scale asset purchases) is the purchase of
long-term securities by a central bank, paid for with reserves. The \emph{term
premium} is the part of a long yield not explained by the expected path of short
rates.
\end{definition}

\begin{record}[Asset purchases, 2008--2014]
% figdata: rates_points
The federal funds rate fell from 4.24~\% in December 2007 to 0.16~\% in December 2008
and stayed there for years. The Federal Reserve then bought medium- and long-term
assets on a large scale. Event studies of the announcements found economically
meaningful and lasting falls in longer-term rates, including on securities it did
not buy, mostly through lower risk and \omterm{def:b3:central-banking:qe}{term premiums} rather than lower expected
short rates. % ledger: b3cb:grrs.lsap
Over the eight announcement days of the first programme, from November 2008 to
November 2009, the ten-year Treasury yield fell by 91 basis points, the yield on
agency mortgage-backed securities by 113 and on ten-year agency debt by 156, while
the two-year yield fell by 34: most of the fall was in long rates, and 71 points
of it in one estimate of the ten-year \omterm{def:b3:central-banking:qe}{term premium}. % ledger: b3cb:grrs.t1
The effects depended on what was bought: mortgage-backed securities purchases in
the first round lowered mortgage yields and corporate credit risk; the second
round, Treasuries only, moved Treasuries most, and other yields mainly through
lower expected policy rates --- a signal, as much as a portfolio effect.
% ledger: b3cb:kvj.channels
\end{record}

\begin{omchart}{Cumulative changes in yields over the eight announcement days of the
Federal Reserve's first asset purchases, November 2008 to November 2009 (basis
points). The term premium is the Kim--Wright model's estimate.}{Gagnon, Raskin,
Remache and Sack (2010), Table~1}
\begin{tikzpicture}
\begin{axis}[omrecord, xbar, bar width=6pt, xmin=-175, xmax=0, ymin=0.4, ymax=7.6,
    y dir=reverse, ytick=data,
    yticklabels from table={figdata/bachelor-3/15-central-banking/out/lsap.csv}{asset},
    y tick label style={font=\footnotesize}, xlabel={Basis points},
    grid=none, xmajorgrids, width=0.6\linewidth, height=0.42\linewidth,
    nodes near coords, every node near coord/.append style={font=\scriptsize, anchor=east}]
  \addplot[fill=omDef!70, draw=omDef] table[x=events, y=pos]
    {figdata/bachelor-3/15-central-banking/out/lsap.csv};
\end{axis}
\end{tikzpicture}
\end{omchart}

An event study sees the announcement and nothing else. Over the whole period from
24 November 2008 to 28 January 2010, as the recovery began, the ten-year Treasury
yield \emph{rose} by 30 basis points; the event days say what the purchases did,
the whole period what everything together did. % ledger: b3cb:grrs.t1
A second method, a time-series regression of the \omterm{def:b3:central-banking:qe}{term premium} on the supply of
long-term debt, put the effect of the 1.725 trillion dollars of announced purchases
at 38 to 82 basis points, against 50 to 100 in the event studies. % ledger: b3cb:grrs.tp

\begin{model}{The portfolio-balance channel of asset purchases}
\assumes Bonds of different maturities are imperfect substitutes: investors who
prefer long bonds pay for them, so a fall in their supply raises their price.
\predicts Purchases lower term premia and long yields, more for the assets bought
and those close to them.
\tested Announcement-day event studies of the first purchases found falls of 91
basis points in the ten-year Treasury yield and 34 in the two-year, mostly in
premia; a time-series model of the \omterm{def:b3:central-banking:qe}{term premium} gave a similar range.
% ledger: b3cb:grrs.t1, b3cb:grrs.tp
\breaks Later rounds, anticipated by markets, moved yields less on the day; part of
the effect works through signals about the policy rate, not through premia; and an
event study sees only the announcement, not the slow effect of holding the bonds.
% ledger: b3cb:kvj.channels
\end{model}

\section{Measuring the effect of a policy shock}\label{sec:b3:central-banking:shocks}

\begin{definition}[Monetary policy shock; local projection]\label{def:b3:central-banking:shock}
\index{monetary policy shock}\index{local projection}
A \emph{monetary policy shock} is the part of a policy decision that markets did not
expect. A \emph{local projection} estimates its effect at each horizon $h$ by a
separate regression of the outcome $h$ periods ahead on the shock and controls.
\end{definition}

Most policy moves respond to the economy, so a regression of output on the policy
rate confuses cause and response. The high-frequency approach measures the change
in interest-rate futures in a narrow window around a scheduled announcement:
within 30 minutes, nothing but the announcement moves rates. After such a surprise
tightening, nominal and real rates rise about one for one several years out, and
expected inflation barely moves --- but forecasts of output growth rise, the
opposite of what a tightening should do, because the announcement also tells
markets what the central bank knows about the economy. % ledger: b3cb:ns.hfi
Identifying a shock is half of the problem; separating the policy news from the
information about the economy it carries is the other.

Before high-frequency surprises were available, the usual shock was a residual: the
part of the policy rate not explained by the current state of the economy and its
recent past. That assumes the central bank sees this quarter's output and inflation
before it moves, and that they do not respond to its move within the quarter. The
listing estimates a \omterm{def:b3:central-banking:shock}{local projection} with such a shock on U.S. quarterly data,
stopping in 2007, before the policy rate reached its floor.

\omcode{code/bachelor-3/15-central-banking/listings/local_projection.py}{1}{30}{A \omterm{def:b3:central-banking:shock}{local projection} of output and prices on a recursively identified rate shock, 1960--2007.}\label{lst:b3:central-banking:lp}

\omoutput{code/bachelor-3/15-central-banking/outputs/local_projection.txt}{Responses in percent, per point of the bill rate; HAC standard errors in parentheses.}\label{out:b3:central-banking:lp}

\begin{omchart}{Estimated responses of U.S. real GDP and of the consumer price level to a
surprise rise of one point in the Treasury-bill rate, 1960--2007, with 90~\% bands
(percent, by quarter after the shock).}{Statsmodels macrodata (U.S. national accounts
and CPI); our computation}
\begin{tikzpicture}
\begin{axis}[omrecord, width=0.62\linewidth, height=0.4\linewidth, xmin=0, xmax=16,
  xtick={0,4,8,12,16}, ymin=-3.2, ymax=0.8, xlabel={Quarters after the shock},
  ylabel={Percent}, legend pos=south west, legend style={font=\footnotesize}]
  \addplot[omDef, thick, mark=*, mark size=1pt] table[x=h, y=y] {figdata/bachelor-3/15-central-banking/out/lp.csv};
  \addplot[omThm, thick, mark=square*, mark size=1pt] table[x=h, y=p] {figdata/bachelor-3/15-central-banking/out/lp.csv};
  \addplot[omDef, thin, dashed, forget plot] table[x=h, y=y_lo] {figdata/bachelor-3/15-central-banking/out/lp.csv};
  \addplot[omDef, thin, dashed, forget plot] table[x=h, y=y_hi] {figdata/bachelor-3/15-central-banking/out/lp.csv};
  \addplot[omThm, thin, dashed, forget plot] table[x=h, y=p_lo] {figdata/bachelor-3/15-central-banking/out/lp.csv};
  \addplot[omThm, thin, dashed, forget plot] table[x=h, y=p_hi] {figdata/bachelor-3/15-central-banking/out/lp.csv};
  \draw[guide] (axis cs:0,0) -- (axis cs:16,0);
  \legend{Real GDP, Price level}
\end{axis}
\end{tikzpicture}
\end{omchart}

% output: local_projection
% figdata: lp_stats
The response at $h = 0$ is zero by construction: the ordering assumes that output
and prices do not react within the quarter. After a surprise rise of one point,
output is 0.72~\% lower after two years and back near its path after four; the price
level barely moves for a year, then falls, by 2.0~\% after four years. Output first,
prices later, is what the models of the Year 2 volume predict. The bands are wide,
and the estimate rests on the ordering: if the bank reacts to news that the data of
the quarter do not contain, part of the ``shock'' is its forecast, the problem the
high-frequency surprises were built to solve.

\section{Financial stability}\label{sec:b3:central-banking:macroprudential}

\begin{definition}[Macroprudential policy]\label{def:b3:central-banking:macroprudential}
\index{macroprudential policy}
\emph{Macroprudential policy} uses regulatory tools --- capital buffers that rise in
booms, caps on loan-to-value or debt-to-income ratios --- to limit risks to the
financial system as a whole, rather than to each bank alone.
\end{definition}

A cap on the loan-to-value ratio binds only for borrowers with small deposits: it
slows credit to the households most likely to default, without raising the policy
rate for everyone. It leaves to monetary policy what only it can do --- set the
price of money --- and to regulation the composition of credit.

\begin{record}[Mortgage limits in Ireland, 2015]
In February 2015 Ireland capped the loan-to-income and loan-to-value ratios of new
residential mortgages, with immediate effect; the limits bound for 43~\% of the
mortgage market. Comparing banks and areas more and less affected, Acharya and
co-authors found that mortgage credit moved from low- to high-income borrowers and
from areas of fast to areas of slow house-price growth, cooling the hottest
markets; the banks most affected took more risk in the securities they held and
in their lending to firms, which the limits did not cover. % ledger: b3cb:abcem.ireland
A difference-in-differences across banks and counties, \emph{grade 2}: the policy
did what it was meant to do to mortgages, and part of the risk moved elsewhere.
\end{record}

\section{2021--2023}\label{sec:b3:central-banking:inflation}

\begin{omchart}{Twelve-month consumer price inflation, January 2019 to December 2025:
United States (CPI-U) and the euro area (HICP).}{U.S. Bureau of Labor Statistics
(CUUR0000SA0, inflation ours); Eurostat (prc\_hicp\_manr)}
\begin{tikzpicture}
\begin{axis}[omtime, width=0.7\linewidth, height=0.4\linewidth, xmin=2018.9, xmax=2026.1,
  ymin=-1, ymax=11.5, ylabel={Percent}, legend pos=north west, legend style={font=\footnotesize}]
  \addplot[omDef, thick, mark=none] table[x=t, y=us] {figdata/bachelor-3/15-central-banking/out/inflation.csv};
  \addplot[omThm, thick, mark=none] table[x=t, y=ea] {figdata/bachelor-3/15-central-banking/out/inflation.csv};
  \draw[guide] (axis cs:2019,2) -- (axis cs:2026,2);
  \legend{United States, Euro area}
\end{axis}
\end{tikzpicture}
\end{omchart}

% figdata: inflation_peaks
U.S. inflation peaked at 9.1~\% in June 2022, the euro area's at 10.6~\% in October
2022; by December 2025 they were 2.7~\% and 2.0~\%. Both central banks kept their policy
rates near or below zero through 2021 and raised them in 2022: the ECB's deposit rate went
from 0.00~\% in July 2022 to 4.00~\% in September 2023. % ledger: b3cb:ecb.rates

\begin{debate}[What caused the inflation of 2021--2023]
Was the surge a supply shock that policy could wait out, or demand that policy
should have met sooner?
\begin{itemize}
\item Bernanke and Blanchard estimated a model of prices, wages and expectations
for the United States: most of the surge came from shocks to prices given wages
--- commodity prices driven by strong demand, and sectoral spikes where demand
met constrained supply; overheated labour markets mattered less at first, but
their effects on wages and inflation are more persistent. % ledger: b3cb:bb.inflation
\emph{Grade 3}: a decomposition that rests on the model's structure.
\item Jordà and co-authors compared the United States with other OECD economies:
U.S. core inflation rose from below 2~\% to above 4~\% in 2021 while the OECD average
rose from about 1~\% to 2.5~\%, and the larger U.S. income transfers may account for
about 3 points of U.S. inflation by the end of 2021. % ledger: b3cb:jlnr.fiscal
\emph{Grade 3}: a cross-country comparison with few countries and many other
differences.
\item Shapiro classified price categories by whether price and quantity moved
together (demand) or in opposite directions (supply): of the 4.8-point excess of
U.S. inflation in April 2022 over its pre-pandemic level, supply explained 2.5
points and demand 1.4. % ledger: b3cb:sha.supply
\emph{Grade 3}: the sign restriction is a plausible but untested identification.
\item Benigno and Eggertsson found the Phillips curve steep when vacancies exceed
the number of unemployed, and attributed most of the surge to an exceptionally
tight labour market. % ledger: b3cb:be.nonlinear
\emph{Grade 3}, with a short sample of tight markets.
\end{itemize}
The studies disagree on the share of demand, but agree that supply and demand
acted together; the sample is short and the shocks simultaneous, which is why
the debate is open.
\end{debate}

\begin{dated}{2026}
The ECB's deposit rate has been 2.50~\% since 16 September 2026. % ledger: b3cb:ecb.rates
The federal funds rate averaged 3.88~\% in October 2026 (Federal Reserve, H.15).
% ledger: b3cb:fred.ff
\end{dated}

\section{Exercises}\label{sec:b3:central-banking:exercises}

\begin{exercise}[$\star$]\label{exo:b3:central-banking:1}
In the opening chart, by how many points of GDP did each balance sheet grow from
2007 to 2021? Which shrank faster afterwards?
\end{exercise}

\begin{exercise}[$\star$]\label{exo:b3:central-banking:2}
A central bank buys 100 billion of government bonds from commercial banks. Write
the changes in its T-account and in the banks' T-account.
\end{exercise}

\begin{exercise}[$\star$]\label{exo:b3:central-banking:3}
In the reserve market of the figure (demand 20 at the lending rate, 60 at the
deposit rate, linear between), the corridor is 2.5--3.5~\%. Find the overnight rate
when reserves are 30, 50 and 70.
\end{exercise}

\begin{exercise}[$\star$]\label{exo:b3:central-banking:4}
What market prices would show that \omterm{def:b3:central-banking:forward-guidance}{forward guidance} is not believed?
\end{exercise}

\begin{exercise}[$\star\star$]\label{exo:b3:central-banking:5}
Short rates are expected to be 1~\%, 2~\% and 3~\% over the next three years, and the
\omterm{def:b3:central-banking:qe}{term premium} on a three-year bond is 0.4 points. Find the three-year yield. If the
central bank's guidance lowers the expected rate in years 2 and 3 by half a point,
what is the new yield?
\end{exercise}

\begin{exercise}[$\star\star$]\label{exo:b3:central-banking:6}
In a \omterm{def:b3:central-banking:corridor}{floor system} with the deposit rate at 2~\%, the central bank raises it to 2.5~\%
without changing reserves. What happens to the overnight rate? And in a corridor
with scarce reserves, if only the deposit rate moves?
\end{exercise}

\begin{exercise}[$\star\star$]\label{exo:b3:central-banking:7}
On three announcement days the ten-year yield fell by 22, 47 and 5 basis points and
the two-year yield by 10, 12 and 4. What share of each fall would you attribute to
expected short rates over the next two years? What does the third day suggest?
\end{exercise}

\begin{exercise}[$\star\star$]\label{exo:b3:central-banking:8}
From the BLS and Eurostat data, find the month and level of peak inflation in the
United States and the euro area, and the first month each fell back below 3~\%.
\end{exercise}

\begin{exercise}[$\star\star$]\label{exo:b3:central-banking:9}
What evidence would contradict the portfolio-balance channel?
\end{exercise}

\begin{exercise}[$\star\star\star$]\label{exo:b3:central-banking:10}
Reserve demand is $R(r) = 60 - 40(r - i_D)/(i_L - i_D)$ for $i_D \le r \le i_L$. With $i_D = 2~\%$, $i_L = 3~\%$
and supply $S$, find $r(S)$ for all $S$, and the supply that puts the rate at the
centre. Why does a floor make the rate insensitive to errors in forecasting
reserve demand?
\end{exercise}

\begin{exercise}[$\star\star\star$]\label{exo:b3:central-banking:11}
A central bank holds 500 of ten-year bonds at 1.5~\% and has 500 of reserves, on
which it pays the policy rate. Find its net interest income at policy rates of 0.5~\%
and 4~\%. Who bears the loss, and does it constrain policy?
\end{exercise}

\begin{exercise}[$\star\star\star$]\label{exo:b3:central-banking:12}
A cap of 80~\% on loan-to-value ratios is introduced. Which buyers does it
constrain, and what should happen to their borrowing and to house prices where
they buy? How would you estimate its effect?
\end{exercise}

\section{Problem: The Bank of Brenland's Balance Sheet}\label{sec:b3:central-banking:problem}

\begin{problem}[{Weekend problem --- The Bank of Brenland's balance sheet}]\label{pb:b3:central-banking:bank}
At the end of 2019 the Bank of Brenland holds government bonds of 60 billion
dalens yielding 3.0~\%, loans to banks of 10 at the policy rate and foreign reserves of
30 in kest bonds yielding 3.0~\%; its liabilities are currency 50, \omterm{def:b3:central-banking:balance-sheet}{bank reserves} 40 and
government deposits 10. Its corridor is half a point either side of the policy
rate, and reserve demand is that of the figure (20 at the lending rate, 60 at the
deposit rate, linear between). The policy rate and inflation come from Brenland's
accounts:
% figdata: bank_years, bank_stats
\begin{center}\small
\pgfplotstabletypeset[col sep=comma, columns={year,policy_rate,infl},
  columns/year/.style={column name={Year}, int detect, column type=l},
  columns/policy_rate/.style={column name={Policy rate (\%)}, fixed, fixed zerofill, precision=2},
  columns/infl/.style={column name={Inflation (\%)}, fixed, fixed zerofill, precision=2},
  every head row/.style={after row=\hline}, every last row/.style={after row=\hline}]
  {figdata/bachelor-3/15-central-banking/out/bank_years.csv}
\end{center}

\medskip\noindent\textbf{Part I --- the T-account.}
\begin{enumerate}
\item Write the Bank's T-account. Check that it balances.
\item The government pays a contractor 4 from its deposit account. What changes?
\item A household withdraws 2 of currency from its bank. What changes?
\item The Bank lends 5 more to banks. What changes, and what happens to the
overnight rate at the 2019 policy rate if the rest is unchanged?
\item Which items does the Bank control, and which does it not?
\item Why are reserves not ``money in circulation''?
\end{enumerate}

\medskip\noindent\textbf{Part II --- from corridor to floor.}
\begin{enumerate}[resume]
\item With reserves of 40 and the 2019 policy rate, find the overnight rate.
\item Find it if reserves fell to 30.
\item Find the Bank's 2019 net interest income, with reserves paid the deposit rate.
\item In 2020 the Bank buys 40 of long bonds yielding 2.0~\% from banks. Write the
changes in its T-account.
\item Find the overnight rate after the purchase, with the corridor around the 2019
policy rate.
\item The Bank moves to a floor: the deposit rate is set equal to the policy rate.
Why does this restore control of the overnight rate?
\end{enumerate}

\medskip\noindent\textbf{Part III --- the purchase programme.}
\begin{enumerate}[resume]
\item Before the purchase, short rates were expected to average 2.7~\% over ten years
and the ten-year yield was 3.5~\%. Find the \omterm{def:b3:central-banking:qe}{term premium}.
\item The purchase lowers the \omterm{def:b3:central-banking:qe}{term premium} by 0.3 points. Find the new yield.
\item By what channel, in the model box, does this happen?
\item Write the Bank's net interest income on the floor as a function of the policy
rate $i$ (in percent).
\item Find it in 2020 and in 2021.
\item Above what policy rate does the Bank make a loss?
\end{enumerate}

\medskip\noindent\textbf{Part IV --- inflation overshoots.}
\begin{enumerate}[resume]
\item The Year 2 volume's rule is $i = 1 + \pi + 0.5(\pi - 2) + 0.5\,\text{gap}$. With a zero gap, find
the rule's rate for 2022.
\item Compare it with the actual 5.42~\%. What could justify the difference?
\item Find the Bank's net interest income in 2022 at the actual and at the rule's
rate, and the latter as a share of GDP (335.8 billion dalens).
\item Does a loss stop the Bank from raising its rate?
\item In 2023 the Bank announces that rates ``will stay high until inflation is
clearly back to target''. What market evidence would show it is believed?
\item State the result: the policy rate implied by the rule and the Bank's net
interest income at that rate.
\end{enumerate}
\end{problem}
